\section{Related Work}

\subsection{IoT-Based Air Quality Monitoring Systems}

The proliferation of low-cost IoT platforms has transformed air quality monitoring from sparse, stationary regulatory networks to dense, real-time sensing infrastructures capable of capturing fine-grained spatiotemporal pollution dynamics \cite{garcia2025advancements}. Contemporary IoT-based indoor air quality (IAQ) systems typically integrate metal-oxide semiconductor (MOS) gas sensors, optical particulate matter detectors, and electrochemical cells with microcontroller units (MCUs) such as ESP32, STM32, or Arduino-based boards, coupled with wireless communication modules (Wi-Fi, BLE, LoRa) for data transmission \cite{bakare2026iot, petrica2026realtime}. Peixe and Marques \cite{peixe2024lowcost} conducted a systematic review identifying 47 IoT-enabled IAQ systems deployed between 2018 and 2024, revealing a trend toward multi-parameter sensing (PM$_{2.5}$, PM$_{10}$, CO$_2$, VOCs, NO$_2$) and cloud-based analytics, yet highlighting persistent challenges in sensor calibration, data quality assurance, and energy management \cite{kim2025machine}.

Energy efficiency has emerged as a critical constraint for battery-powered deployments. Przystupa et al. \cite{przystupa2025ensuring} analyzed power consumption profiles of IoT-based IAQ monitors deployed in residential and industrial settings, demonstrating that wireless transmission accounts for 60-80\% of total node energy expenditure, with instantaneous power spikes reaching 200-300 mW during BLE or Wi-Fi transmissions. Their duty-cycling strategy, which alternates between active sensing (10 s) and deep sleep (50 s), achieved 40\% lifetime extension but did not address data quantization or compression as complementary energy reduction mechanisms. Zaid et al. \cite{zaid2025investigating} deployed a 15-node hyper-local IoT sensor network across urban microenvironments in India, transmitting full-resolution (16-bit) sensor readings every 10 seconds to a cloud backend. While their system successfully captured spatial heterogeneity, the resulting data volume (23 GB/month) necessitated continuous cellular connectivity, limiting scalability to remote or bandwidth-constrained areas.

Bagkis et al. \cite{bagkis2025evolving} reviewed evolving trends in low-cost sensor networks and identified intelligent data management—including adaptive sampling, edge preprocessing, and lossy compression—as a critical enabler for next-generation deployments. However, existing implementations employ heuristic approaches (e.g., transmit only when pollutant changes exceed a threshold) without task-aware optimization tailored to downstream predictive models \cite{cowell2025moving}. Our work addresses this gap by co-designing sensor-side quantization and edge-based TCN inference to optimize the end-to-end energy-accuracy trade-off.

\subsection{Edge AI and Deep Learning for Air Quality Prediction}

The application of deep learning to air quality forecasting has evolved from cloud-centric architectures to edge-deployable models optimized for latency, privacy preservation, and computational efficiency \cite{fici2025sensors, himeur2024edge}. Mazinani et al. \cite{mazinani2025air} recently demonstrated that convolutional and recurrent neural networks quantized to INT8 and INT16 formats can achieve prediction accuracy comparable to FP32 cloud-based models while enabling deployment on resource-constrained microcontrollers (STM32H7, ESP32-S3) with $<$2 MB flash memory. Their work employed post-training quantization (PTQ) of trained models using TensorFlow Lite and validated inference on embedded platforms, reporting sub-100 ms latency for hourly PM$_{2.5}$ forecasting. However, their evaluation assumed availability of high-precision (16-bit) sensor inputs and did not investigate the joint impact of quantizing both sensor data streams and model parameters—a critical omission given that sensor-to-gateway transmission energy often dominates end-to-end system costs \cite{khan2024artificial}.

Davoli et al. \cite{davoli2024edge} proposed an edge computing-oriented Web of Things (WoT) architecture for mobile vehicular air quality monitoring, deploying lightweight CNN models on Raspberry Pi 4 gateways. Their system achieved real-time PM and ozone prediction with $<$1 second latency, but energy analysis focused exclusively on computational overhead (GPU/CPU usage) without accounting for upstream communication costs from sensors to edge nodes. Nguyen et al. \cite{nguyen2025aiot} developed an AIoT-based IAQ prediction system combining metaheuristic feature selection (particle swarm optimization) with hybrid CNN-LSTM models, reporting 8-10\% accuracy improvement over standalone LSTM for 1-hour ahead forecasts on a Vietnamese building dataset. Their approach required extensive hyperparameter tuning (72-hour grid search on Tesla V100 GPUs) and assumed unconstrained computational resources, making it impractical for ultra-low-power edge scenarios.

Gunasekar et al. \cite{gunasekar2025power} introduced the Adaptive Recurrent Temporal Optimized Convolutional Learning (ARTOCL) scheme for power-efficient urban air quality monitoring, combining adaptive sampling triggered by pollutant variability with temporal convolution. ARTOCL reduced transmission frequency by 35-50\% compared to fixed-interval sampling but relied on heuristic variance thresholds (e.g., transmit when $\sigma_{\text{recent}} > 1.5 \sigma_{\text{baseline}}$) rather than principled quantization. Additionally, their energy evaluation lacked explicit modeling of radio power consumption and quantization costs. Yildiz et al. \cite{yildiz2025development} developed a real-time IoT forecasting system using gradient boosting and LSTM models for outdoor air quality in Turkey, but their evaluation prioritized prediction accuracy (R$^2 > 0.92$) over energy constraints, transmitting uncompressed 10-minute averages to cloud servers.

In contrast to these approaches, our work explicitly integrates sensor-side dynamic quantization with hardware-aware TCN design, validated on a large-scale dataset with realistic activity annotations.

\subsection{Quantization Techniques for IoT and Embedded Systems}

Quantization has been extensively studied in the context of model compression for edge deployment, with techniques ranging from uniform linear quantization to learned codebooks, mixed-precision allocation, and knowledge distillation \cite{himeur2024edge}. For air quality applications, quantization serves dual purposes: (i) reducing neural network model size and inference latency for on-device processing, and (ii) compressing sensor data payloads to minimize wireless transmission energy. Rojek et al. \cite{rojek2026impact} analyzed the impact of IoT data analytics and edge computing on energy efficiency in smart environments, emphasizing the role of data preprocessing and lossy compression. Their study identified opportunities for 50-70\% energy savings through adaptive sampling and quantization but provided only generic recommendations without domain-specific designs for environmental time series.

In federated learning and distributed sensing contexts, several works have explored communication-efficient quantization. Khan et al. \cite{khan2024artificial} proposed hybrid LEACH protocols for energy-efficient IoT networks, incorporating gradient quantization to reduce uplink transmission costs in distributed machine learning. Their gradient compression scheme (6-8 bits per parameter update) achieved 5-fold communication reduction with $<$2\% accuracy loss on MNIST and CIFAR-10, demonstrating the viability of aggressive quantization for iterative optimization. However, gradient updates exhibit different statistical properties than raw sensor measurements, limiting direct applicability to environmental monitoring where temporal correlation and bounded ranges can be explicitly exploited.

Cruz et al. \cite{cruz2024improvement} improved an edge-IoT architecture for smart health chronic disease management by quantizing physiological signals (heart rate, SpO$_2$, ECG) before transmission to edge gateways. Their uniform 8-bit quantizer achieved 60\% data reduction with minimal information loss, validated on a 30-patient cohort over 6 months. However, physiological signals exhibit distinct dynamics compared to indoor air pollutants: heart rate variability occurs on subsecond timescales with transient spikes, whereas PM$_{2.5}$ and CO$_2$ evolve slowly (characteristic timescales of minutes to hours) with smooth diurnal patterns \cite{karmakar2024exploring}. This motivates our use of first-order residual encoding and adaptive range scaling tailored to pollutant dynamics.

Recent work on dynamic quantization for IoT intrusion detection \cite{guo2023dynamic} demonstrated that per-layer bit-width adaptation based on activation statistics can reduce model complexity by 40-50\% while preserving detection accuracy. However, their focus on discrete classification tasks (binary intrusion detection) does not address regression objectives or continuous environmental measurements central to air quality forecasting.

\subsection{Temporal Convolutional Networks for Time Series}

Temporal Convolutional Networks (TCNs) have gained prominence as efficient alternatives to recurrent architectures for time series modeling, offering parallel computation, stable gradients, and flexible receptive fields through dilated causal convolutions \cite{bai2018tcn}. Unlike LSTMs and GRUs, which process sequences sequentially and suffer from vanishing gradients over long horizons, TCNs employ stacked 1D convolutional layers with exponentially increasing dilation factors to capture long-range dependencies efficiently. The receptive field of a TCN grows exponentially with depth rather than linearly, enabling modeling of temporal horizons spanning hundreds of timesteps with parameter counts orders of magnitude smaller than recurrent models \cite{lea2017tcn}.

Recent applications of TCNs to environmental monitoring have demonstrated competitive or superior accuracy compared to LSTM baselines. Gunasekar et al. \cite{gunasekar2025power} applied TCNs to urban air quality prediction, reporting 5-8\% RMSE reduction compared to bidirectional LSTMs while achieving 3$\times$ faster inference on ARM Cortex-M7 microcontrollers. Their TCN configuration employed 6 residual blocks with dilations up to 32, capturing 5-hour historical context for 1-hour ahead forecasts. However, their evaluation used synthetic sensor data with idealized noise characteristics, leaving open questions about robustness to real-world quantization artifacts.

Mehta et al. \cite{mehta2023nactcn} proposed NAC-TCN, which augments standard dilated causal convolutions with Dilated Neighborhood Attention layers, achieving state-of-the-art results on high-frequency financial forecasting tasks. While attention mechanisms improve expressiveness, they increase computational cost and parameter counts, potentially limiting edge deployability. Our work employs a standard TCN architecture prioritizing efficiency and quantization-awareness over architectural novelty.

\subsection{DALTON Dataset and Indoor Air Quality Modeling}

The DALTON (Data for Air quality, LocaTion, and Occupancy iN low-income communities) dataset represents a landmark contribution to IAQ research, providing 89.1 million timestamped measurements of seven pollutants (PM$_{2.5}$, PM$_{10}$, CO$_2$, NO$_2$, VOC, ethanol, CO) at 1 Hz resolution from 30 sites across India over approximately 13,646 hours \cite{karmakar2024indoor}. Each DALTON node integrates research-grade sensors (Plantower PMS5003, MQ135 gas sensor, SCD30 CO$_2$ sensor) within a custom 3D-printed enclosure, accompanied by rich metadata including room type (kitchen, bedroom, living room), occupancy levels, and activity labels manually annotated for 42\% of the dataset (e.g., cooking, cleaning, smoking, sleeping) \cite{karmakar2024exploring}. Sites span rural villages, suburban towns, and urban slums, capturing socioeconomic and geographic diversity critical for generalization studies.

Karmakar et al. \cite{karmakar2024exploring} analyzed temporal patterns within DALTON, demonstrating that indoor pollutants exhibit strong autocorrelation over 5-10 minute windows (Pearson correlation $\rho > 0.9$ for lag $\leq$ 300 s), punctuated by sharp transients during cooking (PM$_{2.5}$ spikes $>$ 500 $\mu$g/m$^3$ within 60 s) and cleaning activities. They identified diurnal cycles aligned with daily routines, with peak pollution occurring during morning (6-9 AM) and evening (6-10 PM) cooking periods. This temporal structure motivates our choice of 10-minute ahead forecasting horizon and 4-second subsampling interval—sufficient to capture quasi-static trends while reducing transmission frequency.

Recent studies have begun leveraging DALTON for benchmarking predictive models. Camacho-Magri\~n\'an et al. \cite{camacho2025leveraging} trained LSTM and gradient boosting regressors on DALTON subsets for 30-minute ahead IAQ forecasting, achieving RMSE of 6-8 $\mu$g/m$^3$ for PM$_{2.5}$ but without considering quantization or edge deployment constraints. Wang et al. \cite{wang2025indoor} conducted a cross-geographic assessment of IoT-based IAQ monitoring in Montreal and Qatar, highlighting calibration drift as a persistent challenge for low-cost sensors. Tasmurzayev et al. \cite{tasmurzayev2025lowcost} performed a pilot study in Almaty comparing MOS-based IoT sensors against reference instruments (TSI DustTrak), reporting 20-30\% discrepancies attributable to temperature and humidity interference—an orthogonal challenge to quantization that could be addressed through multi-sensor fusion or calibration layers in future work.

To our knowledge, no prior study has systematically designed and validated a joint sensor-side quantization and edge AI scheme tailored to DALTON's characteristics, with explicit energy-accuracy trade-off analysis and generalizability evaluation across diverse deployment contexts.